\section{Stratégie d'identification et modèle empirique}\label{sec:modele}

La stratégie d'identification repose sur un fait institutionnel : la France a parcouru la même marge
fiscale, l'imposition des dividendes des personnes physiques résidentes, dans les deux sens à cinq
ans d'intervalle. La loi de finances 2013 a ramené les dividendes au barème progressif et alourdi
leur imposition, une hausse que nous estimons sur une fenêtre avant de 2010 à 2012 et une fenêtre
après de 2013 à 2016. La loi de finances 2018, en son article 28, a instauré le prélèvement
forfaitaire unique à 30 \%, une baisse que nous estimons sur une fenêtre avant de 2014 à 2017 et une
fenêtre après de 2018 à 2022. La France est le seul pays à offrir ce couple de chocs symétriques et
rapprochés, ce qui fournit le contraste dont la mesure du cliquet a besoin. Les cinq pays de
contrôle, l'Allemagne, les Pays-Bas, l'Espagne, la Belgique et l'Italie, servent de contre-factuel
sur toute la fenêtre 2010 à 2022 : chacun dispose d'un régime forfaitaire du capital instauré avant
la fenêtre, l'\textit{Abgeltungsteuer} allemande à 26,375 \% en 2009, le \textit{Box 3} néerlandais
en 2001, la base de l'épargne espagnole en 2007, le précompte mobilier belge et la \textit{cedolare
secca} italienne, et aucun ne reproduit l'aller-retour hausse puis baisse de la France. Deux d'entre
eux ont néanmoins ajusté leur taux forfaitaire dans la fenêtre, l'Italie portant la \textit{cedolare
secca} à 26 \% en 2014 et la Belgique son précompte mobilier à 30 \% en 2017, sans parcourir la
marge dans les deux sens ; la permutation placebo, qui classe la France parmi les six pays, et la
robustesse au choix du pool de comparaison traitent cette composition. Le
tableau~\ref{table:identification}, en annexe, récapitule cette variation identifiante. Les dates
des textes de loi sont à confirmer avant diffusion.

L'estimateur principal est une double différence à effets fixes bidirectionnels, dite TWFE pour
\textit{two-way fixed effects}. La France est le seul groupe traité et les cinq pays comparables forment le
groupe de contrôle ; chaque réforme a une date de mise en \oe uvre unique. La spécification compare
donc un seul groupe traité à un ensemble de contrôles autour d'une seule date de traitement, ce qui
en fait un schéma à deux périodes et deux groupes généralisé à plusieurs années. Pour une variable
expliquée $y_{it}$,
\begin{equation}\label{eq:did}
  y_{it} = \beta \, (D_c \cdot P_t) + \theta \, s_{it} + \alpha_i + \gamma_t + \varepsilon_{it}.
\end{equation}
Dans cette équation, $y_{it}$ est la variable expliquée de la société $i$ l'année $t$ ; $D_c$ vaut
un si le pays $c$ est la France, le pays traité, et zéro pour un pays de contrôle ; $P_t$ vaut un à
partir de l'année de la réforme et zéro avant ; $s_{it}$ est le logarithme du total de l'actif ;
$\alpha_i$ est l'effet fixe de société ; $\gamma_t$ l'effet fixe d'année ; et $\varepsilon_{it}$ le
résidu. Le coefficient d'intérêt $\beta$ mesure l'effet sur la France après la réforme, au-delà de
la trajectoire commune des contrôles ; son signe attendu suit les hypothèses, nul pour l'entrée de
capital et les marges de substitution, positif puis plus faible pour les marges de rétention au
retour qu'à l'aller.

Le choix du TWFE appelle une justification, car cet estimateur a fait l'objet de critiques récentes.
Lorsque le traitement est échelonné dans le temps sur plusieurs groupes, la double différence à
effets fixes mélange des comparaisons entre unités traitées tôt et tard et peut pondérer certains
effets négativement \citep{dechaisemartin2020,goodmanbacon2021}. Ces critiques ne s'appliquent pas
ici : un seul groupe est traité et chaque choc a une date unique, de sorte qu'aucune comparaison
interdite entre cohortes de traitement n'intervient et que $\beta$ s'interprète sans ambiguïté comme
la double différence entre la France et ses contrôles avant et après la réforme.

La dynamique est lue par une étude d'événement qui remplace l'interaction unique par des
interactions année par année, suivant la spécification à période de référence omise de
\citet{sunabraham2021},
\begin{equation}\label{eq:event}
  y_{it} = \sum_{j \neq j_0} \tau_j \, \bigl(D_c \cdot \mathbf{1}\{t = j\}\bigr)
           + \theta \, s_{it} + \alpha_i + \gamma_t + \varepsilon_{it},
\end{equation}
où $\tau_j$ est l'écart de la France au pays de contrôle l'année $j$, $\mathbf{1}\{\cdot\}$
l'indicatrice, et $j_0$ l'année de référence omise, 2017 pour le choc de 2018 et 2012 pour le choc
de 2013. Chaque coefficient $\tau_j$ se lit comme l'écart France contre contrôles l'année $j$
relativement à l'année de référence. La platitude des $\tau_j$ avant la réforme teste l'hypothèse de
tendances parallèles ; une rupture après la réforme en signe l'effet dynamique.

L'inférence ne repose pas sur l'erreur-type regroupée, car le nombre de pays traités est de un et le
nombre de groupes de six, trop faible pour une inférence asymptotique fiable. Nous reportons la
statistique de Student à titre indicatif, et l'inférence s'appuie sur une permutation placebo de
l'affectation du traitement à chacun des pays, dont la valeur $p$ est le rang de la France parmi les
six pays. Avec six pays, cette valeur $p$ a un plancher mécanique d'un sixième, soit environ
$0{,}167$, ce que nous assumons explicitement. Au niveau agrégé des ménages, l'effet de chaque choc
est également estimé par un contrôle synthétique \citep{abadie2010,abadie2015}, la France étant
reconstruite comme combinaison convexe des donneurs ajustée sur la période antérieure au choc.
